\begin{abstract}
Heterogeneous cloud workers make synchronous distributed learning vulnerable to
long-tail iteration latency.  Sparse flexible codes reduce redundant work, but
they still leave a systems question open: which recoverable row set becomes
visible first at runtime.  We present \textsc{G-port} (Guarded Prefix-Oriented
Runtime), a worker-service controller whose coded arm keeps the
sparse-flexible code fixed while changing the runtime order in which
recoverable row sets are exposed.  It combines
target-specific row scores, worker speed and delay estimates, deadline-aware
matching, exact online decodability checks, and a fixed guard that enables the
adaptive arm only when predicted prefix savings exceed overhead; otherwise it
falls back to static sparse-flexible placement.  We evaluate \textsc{G-port}
with trace-driven simulation and worker-service
prototypes, including TCP, network-stressed TCP, and a direct three-node K3s
deployment.  Across the worker-service evaluation, \textsc{G-port} improves
barrier latency by up to 24.8\% in deployed K3s stress tests, remains positive
across all main TCP, network-stressed TCP, and K3s scale points, and maintains
bounded recovery under communication, interference, cancellation, and
worker-close stress.  These
results show that row-exposure order is a practical scheduling lever inside
fixed coded-computing deployments.
\end{abstract}
